\documentclass[14pt,a4paper]{extarticle}
\usepackage{gost-style}

\newcommand{\university}{Университет ИТМО}
\newcommand{\faculty}{Школа разработки видеоигр}
\newcommand{\discipline}{Дисциплина: архитектура игровых движков}
\newcommand{\worknumber}{3}
\newcommand{\worktheme}{Менеджер ресурсов и загрузка графических ассетов}
\newcommand{\studentgroup}{J3300}
\newcommand{\studentname}{Лазуренко А.В.}
\newcommand{\teachername}{Павлович А.А.}

\begin{document}

\gosttitlepage
  {\university}
  {\faculty}
  {\discipline}
  {\worknumber}
  {\worktheme}
  {\studentgroup}
  {\studentname}
  {\teachername}

\tableofcontents
\newpage

\section{Цель работы}

Целью практического задания является разработка централизованной системы управления ресурсами игрового
движка. В рамках работы требуется организовать загрузку и кэширование графических ассетов: трехмерных
моделей, текстур, шейдеров и материалов. Система ресурсов должна быть интегрирована с существующей
ECS-архитектурой и рендерингом, чтобы сущности сцены могли ссылаться на внешние файлы моделей и материалов,
а движок автоматически подгружал необходимые данные.

\section{Постановка задачи}

В соответствии с заданием необходимо выполнить следующие пункты:

\begin{enumerate}
  \item подключить внешние библиотеки Assimp, \codeword{stb_image} и \codeword{nlohmann/json};
  \item реализовать обобщенную обертку ресурса и менеджер ресурсов с кэшированием;
  \item создать загрузчики для мешей, текстур, шейдеров и материалов;
  \item расширить слой рендера так, чтобы он мог получать геометрию, текстуры и шейдеры из загруженных ресурсов;
  \item расширить \codeword{MeshRendererComponent} ссылками на загружаемые ресурсы;
  \item адаптировать \codeword{RenderSystem} для отрисовки ECS-сущностей с загруженными 3D-моделями;
  \item продемонстрировать отображение нескольких моделей с разными трансформациями;
  \item выполнить дополнительные задания: сериализацию сцены, асинхронную загрузку, материал как отдельный ресурс и бинарный формат кэша.
\end{enumerate}

\section{Используемые средства}

Проект реализован на языке C++20. Сборка выполняется через CMake и Visual Studio. В качестве графического
API используется DirectX 12. Работа с моделями выполняется через Assimp, загрузка изображений через
\codeword{stb_image}, а чтение и запись JSON через \codeword{nlohmann/json}.

\begin{longtable}{|p{0.32\textwidth}|p{0.58\textwidth}|}
  \hline
  Средство & Назначение \\
  \hline
  C++20 & реализация ядра движка, ECS, менеджера ресурсов и загрузчиков \\
  \hline
  CMake & подключение внешних библиотек, сборка \codeword{UwU_Engine} и \codeword{Sandbox} \\
  \hline
  DirectX 12 & создание GPU-буферов, текстур, pipeline state и выполнение отрисовки \\
  \hline
  HLSL & вершинный и пиксельный шейдеры для отображения цветных и текстурированных объектов \\
  \hline
  Assimp & импорт 3D-моделей формата OBJ и чтение данных материалов из MTL \\
  \hline
  stb\_image & декодирование PNG-текстур в массив RGBA-пикселей \\
  \hline
  nlohmann/json & чтение конфигураций, сцен и material-файлов \\
  \hline
\end{longtable}

\section{Общая архитектура решения}

Система ресурсов размещена в каталоге \codeword{UwU_Engine/src/Engine/Resources}. Основная идея состоит в
том, что загрузка файлов отделена от рендеринга. Загрузчики читают данные с диска и формируют структуры
уровня движка, \codeword{ResourceManager} кэширует результат, а ECS-компоненты хранят только ссылки на
ресурсы и runtime-состояние.

\begin{verbatim}
UwU_Engine/src/Engine/Resources/
  Resource.h
  ResourceTypes.h
  ResourceManager.h
  ResourceManager.cpp
  MeshLoader.h / MeshLoader.cpp
  TextureLoader.h / TextureLoader.cpp
  ShaderLoader.h / ShaderLoader.cpp
  MaterialLoader.h / MaterialLoader.cpp
  BinaryResourceCache.h / BinaryResourceCache.cpp
\end{verbatim}

Интеграция с ECS вынесена в отдельный класс \codeword{RenderResourceBinder}. Он связывает
\codeword{MeshRendererComponent} с системой ресурсов. Благодаря этому \codeword{RenderSystem} не занимается
подробностями загрузки файлов и остается системой отрисовки.

\begin{verbatim}
Scene JSON
   |
   v
MeshRendererComponent
   |
   v
RenderResourceBinder
   |
   v
ResourceManager
   |
   +--> MeshLoader
   +--> TextureLoader
   +--> ShaderLoader
   +--> MaterialLoader
   |
   v
IDrawable / DX12Triangle
\end{verbatim}

\section{Подключение внешних библиотек}

Внешние библиотеки размещены в каталоге \codeword{ExternalLibs}. В CMake-файле движка проверяется наличие
заголовков и библиотечных файлов, после чего добавляются include-директории и линковка.

\begin{lstlisting}[style=cppgost,caption={Подключение nlohmann/json в CMake}]
set(UWU_NLOHMANN_JSON_INCLUDE_DIR
    "${UWU_EXTERNAL_LIBS_DIR}/nlohmann_json/include")

if(EXISTS "${UWU_NLOHMANN_JSON_INCLUDE_DIR}/nlohmann/json.hpp")
    target_include_directories(UwU_Engine PRIVATE
        "${UWU_NLOHMANN_JSON_INCLUDE_DIR}")
else()
    message(FATAL_ERROR "nlohmann_json headers not found")
endif()
\end{lstlisting}

Для Assimp подключаются заголовки, \codeword{.lib}-файл для линковки и \codeword{.dll}-файл для запуска.
DLL копируется рядом с \codeword{Sandbox.exe}, так как Windows ищет динамические библиотеки рядом с
исполняемым файлом.

\begin{lstlisting}[style=cppgost,caption={Подключение Assimp}]
target_include_directories(UwU_Engine PRIVATE "${UWU_ASSIMP_INCLUDE_DIR}")
target_link_libraries(UwU_Engine PRIVATE
    "$<$<CONFIG:Debug>:${UWU_ASSIMP_DEBUG_LIB}>")

add_custom_command(TARGET UwU_Engine POST_BUILD
    COMMAND ${CMAKE_COMMAND} -E copy_if_different
        "${UWU_ASSIMP_DEBUG_DLL}"
        "${OUTPUT_DIR}/Sandbox/")
\end{lstlisting}

\section{Resource и ResourceManager}

Класс \codeword{Resource} представляет загруженный ресурс. Он хранит путь, состояние загрузки, ошибку,
данные конкретного типа и время последнего изменения файла. Для поддержки асинхронной загрузки добавлено
состояние \codeword{Loading}.

\begin{lstlisting}[style=cppgost,caption={Состояния ресурса}]
enum class ResourceState
{
    Unloaded,
    Loading,
    Loaded,
    Failed
};
\end{lstlisting}

Обобщенная обертка ресурса выглядит следующим образом:

\begin{lstlisting}[style=cppgost,caption={Общий принцип Resource}]
template<typename T>
class Resource final : public IResource
{
public:
    explicit Resource(std::string path);

    T& GetData();
    const T& GetData() const;

    void MarkLoading();
    void SetDataLoaded(T data, std::filesystem::file_time_type time);
    void MarkFailed(std::string error);

private:
    T m_data;
    std::filesystem::file_time_type m_lastWriteTime{};
};
\end{lstlisting}

\codeword{ResourceManager} хранит отдельные кэши для разных типов ресурсов. Кэш построен на
\codeword{std::unordered_map}: ключом является нормализованный путь к файлу, значением является
\codeword{std::shared_ptr<Resource<T>>}. Если ресурс уже загружен или загружается, повторный запрос
возвращает тот же объект.

\begin{lstlisting}[style=cppgost,caption={Кэши ResourceManager}]
std::unordered_map<std::string, std::shared_ptr<Resource<MeshAsset>>> m_meshCache;
std::unordered_map<std::string, std::shared_ptr<Resource<TextureAsset>>> m_textureCache;
std::unordered_map<std::string, std::shared_ptr<Resource<ShaderAsset>>> m_shaderCache;
std::unordered_map<std::string, std::shared_ptr<Resource<MaterialAsset>>> m_materialCache;
\end{lstlisting}

Публичный интерфейс построен через шаблонные методы \codeword{Load} и \codeword{LoadAsync}. Внутри они
выбирают нужный конкретный загрузчик по типу ресурса.

\begin{lstlisting}[style=cppgost,caption={Интерфейс загрузки ресурсов}]
template<typename T>
std::shared_ptr<Resource<T>> Load(const std::string& path);

template<typename T>
std::shared_ptr<Resource<T>> LoadAsync(const std::string& path);
\end{lstlisting}

\section{Типы ресурсов}

Основные типы ресурсов описаны в файле \codeword{ResourceTypes.h}. Для модели используется
\codeword{MeshAsset}, содержащий геометрию и данные материала, полученные из файла модели. Для текстуры
используется \codeword{TextureAsset}, для шейдера \codeword{ShaderAsset}, для отдельного material-файла
\codeword{MaterialAsset}.

\begin{lstlisting}[style=cppgost,caption={Основные типы ресурсов}]
struct MeshAsset
{
    MeshData mesh;
    MeshMaterialData material;
};

struct TextureAsset
{
    TextureData texture;
};

struct ShaderAsset
{
    ShaderProgramData shader;
};

struct MaterialAsset
{
    MaterialDesc material;
    bool hasBaseColor = false;
    bool hasShaderPath = false;
    bool hasTexturePath = false;
};
\end{lstlisting}

Такое разделение позволяет отдельно загружать геометрию, изображение, исходный код шейдера и описание
материала, а затем объединять их в \codeword{MeshRendererComponent}.

\section{Загрузчик моделей}

Загрузка моделей реализована в \codeword{MeshLoader}. Для импорта используется Assimp. При чтении файла
применяются флаги триангуляции, генерации нормалей, объединения одинаковых вершин и переворота UV-координат.

\begin{lstlisting}[style=cppgost,caption={Чтение модели через Assimp}]
const aiScene* scene = importer.ReadFile(
    path,
    aiProcess_Triangulate
    | aiProcess_GenSmoothNormals
    | aiProcess_JoinIdenticalVertices
    | aiProcess_ImproveCacheLocality
    | aiProcess_FlipUVs);
\end{lstlisting}

После успешной загрузки вершины переводятся в формат движка \codeword{Vertex}. Для каждой вершины
заполняются позиция, нормаль, UV-координаты и цвет. Индексы формируются по треугольным граням.

\begin{lstlisting}[style=cppgost,caption={Формирование вершины меша}]
Vertex vertex{};
vertex.position[0] = sourceMesh->mVertices[i].x;
vertex.position[1] = sourceMesh->mVertices[i].y;
vertex.position[2] = sourceMesh->mVertices[i].z;

if (sourceMesh->HasNormals())
{
    vertex.normal[0] = sourceMesh->mNormals[i].x;
    vertex.normal[1] = sourceMesh->mNormals[i].y;
    vertex.normal[2] = sourceMesh->mNormals[i].z;
}

if (sourceMesh->HasTextureCoords(0))
{
    vertex.uv[0] = sourceMesh->mTextureCoords[0][i].x;
    vertex.uv[1] = sourceMesh->mTextureCoords[0][i].y;
}
\end{lstlisting}

Также загрузчик читает материал модели. Для OBJ-файла это данные из связанного MTL-файла. Если в материале
указана diffuse-текстура, путь к ней сохраняется в \codeword{MeshMaterialData}.

\begin{lstlisting}[style=cppgost,caption={Получение diffuse-текстуры из материала Assimp}]
aiString diffuseTexturePath;
if (AI_SUCCESS == material->GetTexture(aiTextureType_DIFFUSE, 0, &diffuseTexturePath))
{
    out.material.diffuseTexturePath =
        ResolveReferencedPath(modelPath, diffuseTexturePath.C_Str());
    out.material.hasDiffuseTexture = !out.material.diffuseTexturePath.empty();
}
\end{lstlisting}

\section{Загрузчик текстур}

Загрузка текстур реализована в \codeword{TextureLoader}. Изображение читается через \codeword{stb_image} и
принудительно переводится в формат RGBA. Это упрощает дальнейшую загрузку в DirectX 12, так как drawable
ожидает четыре канала.

\begin{lstlisting}[style=cppgost,caption={Загрузка изображения через stb\_image}]
stbi_uc* pixels = stbi_load(path.c_str(), &width, &height,
    &channels, STBI_rgb_alpha);

const int outputChannels = 4;
const size_t byteCount =
    static_cast<size_t>(width) *
    static_cast<size_t>(height) *
    static_cast<size_t>(outputChannels);

out.texture.width = width;
out.texture.height = height;
out.texture.channels = outputChannels;
out.texture.pixels.assign(pixels, pixels + byteCount);
stbi_image_free(pixels);
\end{lstlisting}

Полученный \codeword{TextureAsset} не является GPU-объектом. Он содержит CPU-данные текстуры. Создание
ресурса DirectX 12 выполняется позднее, когда renderer создает drawable.

\section{Загрузчик шейдеров}

\codeword{ShaderLoader} читает исходный код HLSL-файла в строку. Это позволяет менеджеру ресурсов
кэшировать shader source и передавать его в drawable. Компиляция выполняется в DirectX-слое при создании
конкретного drawable.

\begin{lstlisting}[style=cppgost,caption={Чтение исходного кода шейдера}]
std::ifstream file(path, std::ios::binary);
std::ostringstream stream;
stream << file.rdbuf();

out.shader.sourcePath = path;
out.shader.sourceCode = stream.str();
\end{lstlisting}

В \codeword{MaterialDesc} хранятся entry point и profiles для вершинного и пиксельного шейдера. По
умолчанию используются \codeword{VS}, \codeword{PS}, \codeword{vs_5_0} и \codeword{ps_5_0}.

\section{Material resource}

В рамках самостоятельного задания был добавлен материал как отдельный ресурс. Material-файл хранится в JSON
и описывает базовый цвет, путь к шейдеру, путь к текстуре и настройки компиляции шейдера. Для модели Shiba
создан файл \codeword{Sandbox/src/Assets/Materials/Shiba.material.json}.

\begin{lstlisting}[style=jsongost,caption={Пример material-файла Shiba}]
{
  "color": [ 1.0, 1.0, 1.0, 1.0 ],
  "shaderPath": "../Shaders/Color.hlsl",
  "vertexEntry": "VS",
  "pixelEntry": "PS",
  "vertexProfile": "vs_5_0",
  "pixelProfile": "ps_5_0"
}
\end{lstlisting}

\codeword{MaterialLoader} читает этот JSON через \codeword{nlohmann/json}. Относительные пути в material-файле
разрешаются относительно расположения самого material-файла.

\begin{lstlisting}[style=cppgost,caption={Разрешение относительных путей material-файла}]
std::filesystem::path path(referencedPath);
if (path.is_absolute() || std::filesystem::exists(path))
    return path.lexically_normal().string();

const std::filesystem::path ownerPath(ownerFilePath);
return (ownerPath.parent_path() / path).lexically_normal().string();
\end{lstlisting}

Материал не заменяет данные из MTL-файла модели, а дополняет их. В текущей демонстрации material задает
шейдер и базовый цвет, а путь к diffuse-текстуре берется из MTL-файла модели.

\section{Бинарное кэширование ресурсов}

Для ускорения повторной загрузки реализован простой бинарный формат ресурсов. Он находится в классе
\codeword{BinaryResourceCache}. Бинарный кэш создается рядом с исходным файлом:

\begin{verbatim}
Shiba.obj          -> Shiba.obj.uwumesh
Shiba_base.png     -> Shiba_base.png.uwutex
\end{verbatim}

Перед использованием кэша проверяется его актуальность. Если кэш отсутствует или старее исходного файла,
ресурс загружается обычным способом, после чего кэш пересоздается.

\begin{lstlisting}[style=cppgost,caption={Проверка актуальности бинарного кэша}]
bool BinaryResourceCache::IsFresh(
    const std::string& sourcePath,
    const std::string& cachePath)
{
    if (!std::filesystem::exists(sourcePath) ||
        !std::filesystem::exists(cachePath))
        return false;

    return std::filesystem::last_write_time(cachePath) >=
           std::filesystem::last_write_time(sourcePath);
}
\end{lstlisting}

В \codeword{.uwumesh} сохраняются магическая метка формата, версия, количество вершин, количество индексов,
массив вершин, массив индексов и основные данные материала.

\begin{lstlisting}[style=cppgost,caption={Принцип сохранения меша в бинарный кэш}]
WriteMagic(file, kMeshMagic);
WritePod(file, kVersion);
WritePod(file, vertexCount);
WritePod(file, indexCount);

file.write(reinterpret_cast<const char*>(asset.mesh.vertices.data()),
    asset.mesh.vertices.size() * sizeof(Vertex));
file.write(reinterpret_cast<const char*>(asset.mesh.indices.data()),
    asset.mesh.indices.size() * sizeof(uint32_t));

WritePod(file, asset.material.hasDiffuseColor);
WritePod(file, asset.material.hasDiffuseTexture);
WritePod(file, asset.material.diffuseColor);
WriteString(file, asset.material.diffuseTexturePath);
\end{lstlisting}

Для текстур в \codeword{.uwutex} сохраняются ширина, высота, число каналов и массив пикселей. Таким образом,
при повторной загрузке PNG уже не декодируется через \codeword{stb_image}, а читается как готовый массив
RGBA-данных.

\section{Асинхронная загрузка}

В \codeword{ResourceManager} добавлена асинхронная загрузка через метод \codeword{LoadAsync}. Если ресурс еще
не находится в кэше, создается объект \codeword{Resource}, переводится в состояние \codeword{Loading}, после
чего загрузчик запускается в отдельном потоке.

\begin{lstlisting}[style=cppgost,caption={Идея асинхронной загрузки ресурса}]
resource = std::make_shared<Resource<T>>(normalized);
resource->MarkLoading();
cache[normalized] = resource;

std::thread worker(
    [resource, normalized, asyncLoader = std::move(asyncLoader), typeName]()
    {
        LoadIntoResource(resource, normalized, asyncLoader, typeName);
    });
\end{lstlisting}

Пока ресурс не загружен, \codeword{RenderResourceBinder} может использовать placeholder-геометрию. Для
модели это небольшой серый треугольник. После завершения загрузки drawable сбрасывается и создается заново
уже с настоящей геометрией, текстурой и шейдером.

\begin{lstlisting}[style=cppgost,caption={Placeholder-модель при асинхронной загрузке}]
if (!mesh.mesh.vertices.empty())
    return;

mesh.mesh = MeshFactory::CreateTriangle(
    0.35f,
    Color4{ 0.70f, 0.70f, 0.72f, 1.0f });
mesh.placeholderMeshApplied = true;
ResetDrawable(mesh);
\end{lstlisting}

Потоки загрузки сохраняются в \codeword{ResourceManager}. При очистке менеджер дожидается их завершения,
чтобы не оставить фоновые операции после уничтожения кэша.

\section{Интеграция с ECS}

Компонент \codeword{MeshRendererComponent} был расширен так, чтобы хранить не только готовые данные меша и
материала, но и пути к ресурсам, а также runtime-ссылки на загруженные ресурсы.

\begin{lstlisting}[style=cppgost,caption={Поля MeshRendererComponent для ресурсов}]
MeshData mesh;
MaterialDesc material;
std::string meshResourcePath;
std::string materialResourcePath;

std::shared_ptr<Resource<MeshAsset>> meshResource;
std::shared_ptr<Resource<TextureAsset>> textureResource;
std::shared_ptr<Resource<ShaderAsset>> shaderResource;
std::shared_ptr<Resource<MaterialAsset>> materialResource;
std::unique_ptr<IDrawable> drawable;
\end{lstlisting}

Связка ECS-компонента и менеджера ресурсов вынесена в \codeword{RenderResourceBinder}. Этот класс выполняет
последовательную подготовку компонента:

\begin{lstlisting}[style=cppgost,caption={Последовательность подготовки ресурсов}]
bool Update(EntityId entity, MeshRendererComponent& mesh) const
{
    UpdateMaterial(entity, mesh);
    UpdateMesh(entity, mesh);
    EnsurePlaceholderMesh(mesh);
    UpdateTexture(entity, mesh);
    UpdateShader(entity, mesh);
    return IsShaderReady(mesh);
}
\end{lstlisting}

\codeword{RenderSystem} теперь не содержит подробной логики загрузки ресурсов. Он проверяет видимость,
вызывает \codeword{RenderResourceBinder}, создает drawable и выполняет отрисовку.

\section{Рендеринг загруженных моделей}

На уровне DirectX 12 общий код геометрии вынесен из прежнего класса \codeword{DX12Triangle} в
\codeword{DX12MeshGeometry}. Этот класс создает vertex buffer и index buffer для произвольного
\codeword{MeshData}. Благодаря этому drawable больше не ограничен одним заранее заданным треугольником.

\begin{lstlisting}[style=cppgost,caption={Использование общей геометрии в drawable}]
bool DX12Triangle::BuildGeometry(ID3D12Device* device,
    const DrawableDesc& desc)
{
    return m_geometry.Create(device, desc.mesh);
}
\end{lstlisting}

Текстура создается в \codeword{DX12Triangle} из данных \codeword{MaterialDesc}. Если пиксели не переданы,
используется белая текстура 1x1. Это позволяет отрисовывать как текстурированные, так и нетекстурированные
объекты одним путем.

\begin{lstlisting}[style=cppgost,caption={Fallback-текстура}]
std::vector<uint8_t> fallbackWhite = { 255, 255, 255, 255 };
const uint8_t* pixels = fallbackWhite.data();
int width = 1;
int height = 1;

if (!desc.material.texturePixels.empty())
{
    pixels = desc.material.texturePixels.data();
    width = desc.material.textureWidth;
    height = desc.material.textureHeight;
}
\end{lstlisting}

Шейдер \codeword{Color.hlsl} использует позицию, цвет вершины и UV-координаты. В пиксельном шейдере цвет
текстуры умножается на цвет вершины.

\begin{lstlisting}[style=cppgost,caption={Фрагмент пиксельного шейдера}]
float4 PS(VertexOut pin) : SV_Target
{
    float4 diffuse = gDiffuseTexture.Sample(gDiffuseSampler, pin.TexCoord);
    return diffuse * float4(pin.Color, 1.0f);
}
\end{lstlisting}

\section{Сериализация сцены}

Сцена хранится в JSON-файле \codeword{Sandbox/src/Assets/Scenes/scene_test.json}. Для моделей в
\codeword{meshRenderer} указываются путь к мешу, путь к material-файлу, путь к шейдеру или текстуре при
необходимости, а также базовый цвет.

\begin{lstlisting}[style=jsongost,caption={Фрагмент сцены с моделью Shiba}]
{
  "id": 6,
  "tag": { "name": "Shiba model" },
  "transform": {
    "x": 0.85,
    "y": 0.18,
    "z": 0.0,
    "rotationY": 1.5708,
    "scaleX": 1.25,
    "scaleY": 1.25,
    "scaleZ": 1.25
  },
  "meshRenderer": {
    "primitive": "CustomMesh",
    "meshPath": "Assets\\Models\\Shiba\\Shiba.obj",
    "materialPath": "Assets\\Materials\\Shiba.material.json",
    "shaderPath": "",
    "texturePath": "",
    "color": [ 1.0, 1.0, 1.0, 1.0 ]
  }
}
\end{lstlisting}

В сцене добавлены две сущности с моделью Shiba. Первая модель вращается по оси Y в игровом цикле, вторая
является статичной и имеет другой масштаб. Это демонстрирует повторное использование одного и того же
ресурса модели разными ECS-сущностями.

\section{Демонстрационная сцена}

Демонстрационная сцена содержит:

\begin{enumerate}
  \item несколько треугольников из предыдущих практических работ;
  \item основную управляемую сущность;
  \item камеру с перемещением и масштабированием колесом мыши;
  \item вращающуюся модель Shiba;
  \item статичную модель Shiba с другим масштабом;
  \item material-файл, общий для моделей Shiba;
  \item текстуру \codeword{Shiba_base.png}, путь к которой берется из MTL-файла модели.
\end{enumerate}

\placeholderimage{Окно приложения с двумя моделями Shiba}{Место для скриншота демонстрационной сцены}

\placeholderimage{Фрагмент логов загрузки ресурсов}{Место для скриншота логов ResourceManager, MeshLoader и TextureLoader}

\section{Трудности и их решение}

В процессе выполнения работы возникло несколько практических проблем.

Первая проблема была связана с путями к ассетам. При запуске из Visual Studio и при запуске из
\codeword{Sandbox.exe} текущая рабочая директория может отличаться. Для решения CMake копирует папку
\codeword{Assets} рядом с исполняемым файлом, а загрузчики нормализуют относительные пути.

Вторая проблема возникла при использовании бинарного кэша. В \codeword{.uwumesh} мог быть сохранен путь к
текстуре, актуальный для запуска из исходной папки. При запуске из \codeword{bin} такой путь становился
невалидным. Поэтому при чтении mesh-кэша путь к diffuse-текстуре дополнительно пере-привязывается к текущему
расположению \codeword{.obj}-файла и папке \codeword{Assets/Textures}.

Третья проблема была связана с отображением 3D-моделей при движении камеры. Перспективная проекция создавала
визуальное изменение формы модели при панорамировании. Для текущей демонстрации была выбрана ортографическая
камера, чтобы движение камеры не изменяло пропорции и видимый угол статичных объектов.

\section{Результаты}

В результате практической работы реализована система ресурсов игрового движка:

\begin{enumerate}
  \item подключены Assimp, \codeword{stb_image} и \codeword{nlohmann/json};
  \item создан \codeword{ResourceManager} с кэшированием ресурсов;
  \item реализованы загрузчики мешей, текстур, шейдеров и материалов;
  \item \codeword{MeshRendererComponent} расширен ссылками на ресурсы;
  \item \codeword{RenderSystem} адаптирован для отрисовки загруженных моделей;
  \item добавлен \codeword{RenderResourceBinder}, отделяющий подготовку ресурсов от логики рендера;
  \item реализована асинхронная загрузка ресурсов;
  \item добавлен placeholder на время загрузки модели;
  \item реализован бинарный формат кэша \codeword{.uwumesh} и \codeword{.uwutex};
  \item создан material resource и material-файл для модели Shiba;
  \item в сцене отображаются две 3D-модели Shiba с разными параметрами трансформации.
\end{enumerate}


\section{Заключение}

В ходе практического занятия была разработана и интегрирована система управления ресурсами игрового движка.
Реализация позволила перейти от ручной отрисовки простых примитивов к загрузке внешних 3D-моделей и текстур.
Архитектура стала более расширяемой: загрузка ресурсов отделена от рендера, материалы вынесены в отдельные
файлы, а ECS-сцена хранит ссылки на ассеты.

Полученная система может быть расширена в следующих работах. В качестве дальнейших направлений можно
выделить поддержку нескольких материалов внутри одной модели, normal map, skeletal animation, полноценный
asset database и горячую замену ресурсов в runtime.

\section{Приложения}

Основные файлы, измененные или созданные при выполнении работы:

\begin{verbatim}
UwU_Engine/src/Engine/Resources/Resource.h
UwU_Engine/src/Engine/Resources/ResourceTypes.h
UwU_Engine/src/Engine/Resources/ResourceManager.h
UwU_Engine/src/Engine/Resources/ResourceManager.cpp
UwU_Engine/src/Engine/Resources/MeshLoader.cpp
UwU_Engine/src/Engine/Resources/TextureLoader.cpp
UwU_Engine/src/Engine/Resources/ShaderLoader.cpp
UwU_Engine/src/Engine/Resources/MaterialLoader.cpp
UwU_Engine/src/Engine/Resources/BinaryResourceCache.cpp
UwU_Engine/src/Engine/ECS/RenderResourceBinder.h
UwU_Engine/src/Engine/ECS/Components.h
UwU_Engine/src/Engine/ECS/RenderSystem.h
Sandbox/src/Assets/Scenes/scene_test.json
Sandbox/src/Assets/Materials/Shiba.material.json
\end{verbatim}

Ссылка на репозиторий:

\begin{verbatim}
https://github.com/<username>/<repository>
\end{verbatim}

\end{document}
